% Transcription — fonds Alexandre Grothendieck, shelfmark 77, batch 6 (pages 101-120).
% Established from the facsimile archives/batches/77/batch-06.pdf (a mirror of
% https://grothendieck.umontpellier.fr/77.pdf), per the skill
% transcribe-grothendieck.
%
% Pass: Opus 5.5 (claude-opus-5-5), 2026-10-03 — first pass, unchecked against the pages by a human.
%
\documentclass[11pt,a4paper]{article}
\input{../preamble/grothendieck}

\folder{77}
\batch{6}
\pages{101}{120}
\dating{[à partir de 1977]}
\watermark{Édition de démonstration}
\foldertitle{Quadriques affines, extensions quadratiques : notes manuscrites (s.d.).}

\begin{document}

\page{101}
\note{La page continue un calcul commencé avant ce lot (les notations $\tilde{U}$, $\Sigma_{-4}$, $\varphi$, $[x,y]$ y sont déjà fixées). La moitié supérieure, jusqu'à la flèche $g$, est barrée de longs traits obliques ; elle se lit néanmoins et on la donne entière sous \texttt{struck}.}

\textbf{Car. \ill{}}

\struck{$q([x,y]) = \varphi(x,y)^{2}$.}
\struck{Donc si $x + \mathbf{i}y \in \Gamma\tilde{U}$, donc $q(x) = q(y)$, $\varphi(x,y) = 0$} \add{\struck{$\in \Gamma\,\mathbb{G}_{m}$}}\struck{, alors $q([x,y]) = 0$. Si \ill{} remplace $x + \mathbf{i}y$ par}
\[
  \text{\struck{$(a + \mathbf{i}b)(x + \mathbf{i}y) = \underbrace{ax - by}_{x'} + \mathbf{i}(\underbrace{bx + ay}_{y'})$,}}
\]
\struck{on trouve $x' \wedge y'$ $= (a^{2} + b^{2})\, x \wedge y$, $q(x') = q(y')$ $= (a^{2} + b^{2})\, q(x)$, donc $x' \wedge y' / q(x')$ $= x \wedge y / q(x)$. On trouve}
\[
  \text{\struck{$x + \mathbf{i}y \in \Gamma\tilde{U} \Longrightarrow [x,y]/q(x) \in \Sigma_{-4}$}}
\]
\struck{donc}
\[
  \text{\struck{$\tilde{g} : \tilde{U} \longrightarrow \Sigma_{-4}$}}
\]
\struck{passe au quotient en}
\[
  g : \text{\struck{$\tilde{X}$}}\ \longrightarrow \Sigma_{-4}
\]
(\uncertain{indépendamment} de la \ill{}).
\struck{S\ill{}, $[x,y]/q(x) = [x',y']/q^{2}(x')$ ; \ill{}}
\struck{(quitte à localiser, et à remplacer $x'$ par $\lambda x'$ \ill{}, $q(x) = q(x')$, $[x,y] = [x',y']$. S\ill{} que $g \vert S$ \ill{} \ill{} car.}\note{Ce bloc encadré et barré est peu lisible ; les fragments donnés ici sont ce qui se lit sous les ratures.}
\[
  [x,y] = \varphi(x,y)\, H \quad \text{où} \quad
  \begin{cases}
    H \text{ base du centre de } \tilde{\mathfrak{g}} \\
    q(H) = 1
  \end{cases}
\]
donc
\[
  q([x,y]) = \varphi(x,y)^{2}.
\]
Si $\varphi(x,y) = 0$ (\ill{} $x + \mathbf{i}y \in \Gamma\tilde{X}^{S[i]} \supset \Gamma(\tilde{U})$), \ill{} $[x,y] = 0$. Dans car.~2 l'involution se spécialise en l'involution constante du sous-schéma $0$ : rien à faire…!

\page{102}
\textbf{(1')} Soit $S$ schéma, $E$ fibré affine de rg~3 sur $S$, $\Sigma \subset E$ une quadrique affine \add{(projectiv\supplied{emen}t \add{lisse} \struck{régulière})}. \ill{} il revient au $=$ de \struck{\ill{}} \add{($= \hat{E}$)} se donner une partie d'un fibré projectif $P$ de rang~3, d'une quadrique \add{lisse} \add{($= \hat{\Sigma}$)} $Q \subset P$, et d'un hyperplan $H$ de $P$ --- on prend $E = P \setminus H$, $\Sigma = Q \cap E = Q \setminus Q \cap H$. \struck{\ill{}} [\uncertain{Si on veut}, il revient encore au $=$ de se donner $Q$ comme schéma (projectif lisse) sur $S$ qui soit une forme de $\mathbb{P}^{1} \times \mathbb{P}^{1}$, ce qui détermine $P$ canoniquement, et de prendre un diviseur \add{$X$} de type $(1,1)$ sur $Q$ --- lequel est défini par un $H$ bien déterminé dans $P$…]
\note{Le signe $=$ isolé (« revient au $=$ ») est son abréviation de « même ».}

On suppose \struck{\ill{}} que fibre par fibre $\Sigma$ n'est pas un \uncertain{paraboloïde}, i.e.\ que $H$ n'est pas tangent à $Q$, ou encore que $X = H \cap Q = \hat{\Sigma} \setminus \Sigma$ est lisse, i.e.\ un \ill{} projectif sur $S$. On peut dire aussi que la donnée de $(E, \Sigma)$ équivaut à \uncertain{celle d'une} ext.\ $\mathcal{E}$ de $\mathcal{O}_{S}$ par

\page{103}
un fibré vectoriel de rg~3 $V$ (fibré des translations de $E$)
\[
  0 \longrightarrow V \longrightarrow \mathcal{E} \longrightarrow \mathcal{O}_{S} \longrightarrow 1
\]
\note{Le dernier terme est écrit comme un trait vertical, qui se lit $1$ ; on attendrait $0$.}
et « une » forme quadratique \add{lisse} $Q$ (sur $\mathcal{E}$) définie loc\supplied{alemen}t \uncertain{à mult.\ près par une unité}. [\ill{} \emph{signifie} que $\Sigma$ est définie par une fonction polynomiale de degré $\leq 2$, et qui est projectiv\supplied{emen}t lisse…]

La condition que $\Sigma$ ne soit pas un \uncertain{paraboloïde} signifie que $q \vert V$ est \struck{\ill{} dégénérée} lisse. Si les car.\ résid.\ sont $\neq 2$, cela signifie \add{aussi} que
\[
  \mathcal{E} = V \oplus V^{\perp},
\]
donc $E$ se trouve muni \add{\uncertain{parallèlement à l'idéal} : $V$} d'une origine $o$ (le pôle de $H = P(V)$) \add{(\uncertain{le centre} de $Q$)} et $\Sigma$ est \uncertain{stable par la} symétrie p.r.\ à cette origine --- $o$ s'appelle le \emph{centre} de la quadrique affine. Ce \struck{\ill{}} \add{section} de $E$ sur $S$) est rendu fonctoriel par $\Sigma$, \struck{la} la condition quadratique peut être normalisée par la condition \uncertain{qu'elle} vaut~$1$ \uncertain{en} $o$, ce qui la détermine de

\page{104}
façon unique et assure son existence. On \struck{\ill{}} trouve donc que sur $E \;(= V \times 1)$ \struck{\ill{}} \add{$\subset \mathcal{E}$} la fonction \uncertain{induite} par $Q$ est de la forme $q(x) + 1$ où $q$ est la forme quadratique \add{lisse} induite sur $V$ par $Q$. Donc la donnée de $\Sigma$ dans $E$ équivaut (au \ill{} des possibilités) à celle de

a) Section $\sigma$ de $E$ (le \emph{centre} de $\Sigma$)

b) forme quadratique \add{lisse} sur $V$ (fibré vectoriel des translations)

On associe à $\sigma$, $q$ la quadrique
\[
  \Sigma = \check{\mathbb{V}}(q+1) = \{x \in \Gamma(E, V) \mid q(x) = -1\}
\]
(en identifiant $E$ à $V$ grâce à $\sigma$).

\note{Le paragraphe suivant est marqué en marge gauche d'un double trait vertical, à l'encre noire.}
Ces développements sont valables \ill{} \ill{} (\ill{} de $E$) au \ill{} \uncertain{premier sens}, \ill{} \add{(\uncertain{comme} \ill{} \ill{})}, \ill{} \ill{} \ill{} \ill{}, ce \ill{} des car.\ résid.\ $\neq 2$.

Considérons
\[
  \delta'(q^{\ast}) \in (\det \check{V})^{\otimes 2}
  \qquad \text{base de } (\det \check{V})^{\otimes 2},
\]
\note{L'objet écrit après « Considérons », une sorte de $\delta$ surmonté d'un trait, est surchargé ; l'exposant de $q$ est peu lisible.}

\page{105}
et considérons le \ill{} de
\[
  \check{\mathbb{V}}(\det \check{V}) = \mathbb{V}(\det(\check{V})) = \mathbb{V}(\det(V)^{-1})
\]
défini par l'équation
\[
  \omega^{\otimes(-2)} = \delta'(q).
\]
C'est un torseur sous $\mu_{2,S} \simeq (\mathbb{Z}/2\mathbb{Z})_{S}$, i.e.\ un rev\supplied{êtement} étale de degré~2 de $S$, \marginal{(car.\ résid.\ $\neq 2$ !)} noté $\Omega$. [Dans le cas où $q$ est spécial orthogonal, i.e.\ muni de $\omega_{0}$ telle que $\omega_{0}^{\otimes -2} = -\delta'(q)$ --- attention au signe ! --- les $\omega$ correspondent aux $\lambda\omega_{0}$ avec $\lambda^{2} = -1$, dans $\Omega = S[\mathbf{i}] \overset{\text{déf}}{=} \operatorname{Spec} \mathcal{O}_{S}[T]/(T^{2}+1)$.]
\note{Le signe devant $\delta'(q)$ est doublé d'un second trait, sans doute pour le souligner ; le $i$ de $S[i]$ est souligné, comme ailleurs dans le lot : on le rend par $\mathbf{i}$.}

Considérons la clôture projective $\hat{\Sigma}$ ($= Q$) de $\Sigma$, et \struck{\ill{}} sa \ill{} à l'infini, \struck{\ill{}} \add{conique} projective $X$ \add{($\subset H = \check{\mathbb{P}}(V)$, hyperplan à l'infini de $\hat{E}$)} \uncertain{qui est aussi} \ill{} de $\hat{E} = \check{\mathbb{P}}(\mathcal{E})$, définie par la forme $q$. On voit que, si $\sigma$ désigne la \struck{\ill{}} symétrie p.r.\ à l'origine $\sigma$ de $E$, \ill{} \ill{} $\hat{E}$), on a
\[
  \hat{E}^{\sigma} = \sigma(S) \sqcup H, \quad \text{(}H = \mathbb{P}(V)\ \text{hyperplan à l'infini de } \hat{E}\text{)}
\]
\[
  \hat{\Sigma}^{\sigma} = H \cap \hat{\Sigma} = X.
\]
\note{Le membre « $= \sigma(S) \sqcup H$, … » est ajouté à l'encre noire au-dessus de la parenthèse sur $H$.}

\page{106}
\struck{Notre propos est de définir un isomorphisme}

Considérons \struck{$X$} d'autre part
\[
  \mathbf{Hom}_{S}(\Omega, X) \overset{\text{déf}}{=} X^{\Omega} ;
\]
l'automorphisme can.\ de $\Omega$ définit un automorphisme de $X^{\Omega}$, noté \add{également} $\sigma$, \struck{donc} et on a évidemment un iso.\ can.
\[
  (X^{\Omega})^{\sigma} \simeq X.
\]
Ceci posé, notre propos est de définir un iso canonique
\[
  X^{\Omega} \xrightarrow{\ \sim\ } \hat{\Sigma}
\]
compatible avec $\sigma$, et rendant commutatif le diagramme

\begin{tikzcd}
  & X \arrow[dl, hook'] \arrow[dr, hook, "i"] & \\
  X^{\Omega} \arrow[rr, "\sim"] & & \hat{\Sigma}
\end{tikzcd}

donc induisant un iso
\[
  \Sigma \simeq X^{\Omega} \setminus X,
\]
compatible avec $\sigma$.

\page{107}
Pour définir un iso
\[
  X^{\Omega} \xrightarrow{\ \sim\ } \hat{\Sigma}
\]
il revient au $=$, par descente, de définir un iso
\[
  (X^{\Omega})_{\Omega} \xrightarrow{\ \sim\ } (\hat{\Sigma})_{\Omega}
\]
compatible avec les $\Omega$-autom.\ des deux schémas, ce qui revient au $=$ [comme
\[
  (X^{\Omega})_{\Omega} \simeq (X_{\Omega})^{\Omega_{\Omega}}, \quad
  \Omega_{\Omega} = \Omega \times_{S} \Omega = \Omega \sqcup \Omega,
  \quad \simeq X_{\Omega} \times X_{\Omega},
\]
avec l'action de descente (provenant de l'action de $\Omega$ à $X^{\Omega}$) se transformant en l'action composée de la symétrie des 2 facteurs et de l'automorphisme induit $\sigma_{\Omega}$ de $\Omega$,] il revient au $=$ de trouver un isom.
\[
  X_{\Omega} \times X_{\Omega} \xrightarrow{\ \sim\ } \hat{\Sigma}_{\Omega} \quad (= \widehat{\Sigma_{\Omega}})
\]
qui transforme $\sigma_{\Omega}^{(\ast)} \times s$ en $\sigma_{\Omega}^{(\hat{\Sigma}_{\Omega})}$.
\note{L'exposant du premier $\sigma_{\Omega}$ est un mot abrégé, illisible (« \ill{} ») ; $s$ est la symétrie des deux facteurs.}
Cela signifie aussi : pour $\omega \in \Gamma(S', (\det V)^{\otimes 2})$ \ill{} ($S'$ \uncertain{quelc.}\ sur $S$) tel que
\[
  \delta'(q) = \omega^{\otimes 2},
\]
trouver un iso
\[
  (\ast) \qquad f_{\omega} : X \times X \longrightarrow \hat{\Sigma}
\]
\struck{\ill{}} satisfaisant les conditions
\note{Dans $\delta'(q)$, le $q$ est récrit à l'encre noire par-dessus un symbole biffé. L'exposant de $\det V$ dans $\Gamma(S', \ldots)$ est peu sûr.}

\page{108}
a) compatibilité au changement de base \struck{$f(-\omega, x, y) =$} \struck{\ill{}}

b) $f_{-\omega}(x,y) = f_{\omega}(y,x)$

c) $f_{\omega}(y,x) = \sigma f_{\omega}(x,y)$

d) $f_{\omega}(x,x) = i(x)$

\note{Les lettres b) et c) sont surchargées l'une sur l'autre ; en c), le membre de gauche est d'abord écrit $f_{\omega}(x,y)$ puis corrigé.}
\struck{\ill{}} Notons que \add{donc b) et c) \uncertain{impliquent ensemble}}
\[
  f_{-\omega}(x,y) = f_{\omega}(y,x) = \sigma f_{\omega}(x,y),
\]
\[
  \delta'(-q) = -\delta'(q) = -\omega^{\otimes 2},
\]
donc \struck{il y a un iso} on peut trouver, grâce à $\omega$, une \uncertain{algèbre} de \struck{Azumaya} \ill{} $M$ et un iso
\[
  E \overset{\varphi}{\simeq} \gamma M
\]
\marginal{attention au signe}
transformant $-q$ en $q_{M} = -\det \vert \gamma M$ et $\omega$ en $\omega_{M}$, donc transformant \add{$(M, \varphi)$ \ill{} \ill{} : iso\supplied{morphisme} unique près par $q, \omega$} \ill{} $X$ s'interprète comme la variété projective associée au cône des \add{\ill{}} nilpotents de $M$, définis par les équations
\[
  \operatorname{Tr} u = \det u = 0.
\]
Quant à $\Sigma$, il correspond à l'équation $\det u = -1$ en plus de $\operatorname{Tr} u = 0$, qui correspond au schéma des \emph{réflexions} associées à $M$. L'iso.\ induit par $f_{\omega}$ donne un iso
\[
  (X \times X) \setminus X \simeq \Sigma
\]
s'interprétant comme une bijection (fonctorielle

\page{109}
pour le changement de base)
ens.\ des couples $(u, v)$ d'él.\ de $M$ \struck{\ill{}}
(i.e.\ \ill{} du tr.\ et dét.\ nuls), \add{(\ill{} \uncertain{scalaires} \ill{})} \uncertain{distincts} fibre par fibre
\[
  \simeq \text{ens.\ des réflexions de } M \quad (\operatorname{tr} = 0,\ \det = -1).
\]
\struck{Mais les réflexions correspondent \ill{}}

Si on veut passer aux clôtures projectives, on note que (comme car.\ rés.\ $\neq 2$) \ill{}
\[
  M \simeq \gamma M \oplus \mathcal{O}_{S},
\]
\note{Sous $\gamma M$, un signe $\wr$ relie à $\tilde{V}$ : $\gamma M \simeq \tilde{V}$.}
\[
  \hat{\Sigma}_{M} \subset \check{\mathbb{P}}(M)
\]
défini par l'équation \add{\uncertain{dans} $\check{\mathbb{V}}(\gamma M \oplus \mathcal{O}_{S})$}
\[
  \det(u) + t^{2} = 0 \qquad \text{i.e.} \qquad \det(u + t\,\mathrm{id}) = 0 ;
\]
\note{Devant $t^{2}$, un coefficient est biffé.}
si on veut, on pose $v = u + t\,\mathrm{id}$,

i.e.\ \struck{c'est} \add{c'est la} quadrique projective définie par $\det$ de $M$, $\det v = 0$. On sait qu'on a un iso.\ canonique \add{\ill{} celle-ci et}
\[
  F_{M} : X_{M} \times X_{M} \simeq \hat{\Sigma}_{M}, \qquad
  \hat{\Sigma} \cap \mathbb{P}(M) \simeq X
\]
\struck{avec} (qui sur $X_{M} \times X_{M}$) induit l'application diagonale. On sait de plus que \struck{$F_{M}(y,x)$} $F_{M}$ transforme la symétrie de $X_{M} \times X_{M}$ en l'\uncertain{effet} sur $\hat{\Sigma}_{M}$ de l'antiinvolution
\[
  v \longmapsto v' = \operatorname{tr}(v).1 - v
\]
(qui conserve le déterminant…), ce qui \add{($u \in \Gamma\,\gamma M$, $t \in \Gamma(\mathcal{O}_{S})$)} pour
\[
  v = u + t1
\]
s'écrit
\[
  v' = 2t.1 - (u + t.1) = -u + t1,
\]
soit
\[
  -v' = u - t.1
\]

\page{110}
donc l'automorphisme de $\hat{\Sigma}_{M}$ induit par $v \mapsto v'$ est aussi l'antipodisme, induit par la réflexion
\[
  \underset{v}{\underbrace{u + t.1}} \longmapsto \underset{-v'}{\underbrace{u - t.1}} .
\]
Ceci \struck{\uncertain{par transport}} échange \uncertain{du dessin} notre iso.\ $(\ast)$, satisfaisant a) c) d) ; il reste à vérifier b), ce qui revient au $=$, la formule
\[
  f_{-\omega}(x,y) = \sigma f_{\omega}(x,y),
\]
ce qui résulte du fait que $(\omega, x, y) \mapsto f_{\omega}(x,y)$, $X \times X \times \Omega \longrightarrow \hat{\Sigma}$, \uncertain{étant} défini fonctoriellement en $(E, q, \omega)$, commute à l'action de l'iso
\[
  -\mathrm{id}_{E} : (E, q, \omega) \longrightarrow (E, q, -\omega).
\]
\note{Sous les deux membres de la flèche $u + t.1 \mapsto u - t.1$, il écrit $\overset{\shortparallel}{v}$ et $\overset{\shortparallel}{-v'}$ ; on les rend par des accolades.}

Il reste : faire la jonction avec les développements du §1, relatifs au cas où on se donne, au $=$ temps que $q$ \add{sur $V$}, définissant la quadrique affine d'équation $q(x) + 1 = 0$, un $\omega_{0} \in \Gamma \det \check{E}$ satisfaisant (non pas $\delta'(q) = \omega_{0}^{\otimes 2}$ mais…)
\[
  \delta'(q) = -\omega_{0}^{\otimes 2} \qquad \text{i.e.} \qquad \delta'(-q) = \omega_{0}^{\otimes(-2)}
\]
\note{Dans la dernière formule, l'exposant est d'abord $\otimes 2$, puis $\otimes(-2)$ ; $\omega_{0}$ y est surchargé.}

\page{111}
Dans ce cas, montrons que l'iso
\[
  X^{\Omega} \setminus X \overset{f_{q,\omega}}{\simeq} \Sigma \qquad (\text{où } \Omega \simeq S[\mathbf{i}])
\]
induit par l'iso $X^{\Omega} \simeq \hat{\Sigma}$ est donné par la formule du § précédent, \struck{\ill{}} soit $g_{q}$.

\note{Le bloc qui suit, jusqu'à $\theta : \gamma M \to \gamma M'$, est encadré à gauche et barré de deux longues diagonales ; il reste lisible et on le donne sous \texttt{struck}.}
\struck{Ors $(E, q, \omega_{0})$ de la forme $(\gamma M, q_{M}, \omega_{M})$, donc $\Sigma$ défini par $\det u = 1$, $\operatorname{tr} u = 0$ [i.e.\ c'est le schéma des opérations $u \in M$ \add{de H.C.} d'équation \struck{\ill{}} $u^{2} + 1 = 0$, i.e.\ \struck{qui} \struck{définissant} schém.\ des \struck{structures de} inclusions $\mathcal{O}_{S}[\mathbf{i}] \hookrightarrow M$ …] La donnée de $\omega$, \add{\uncertain{revient} : \ill{} \uncertain{dessus}} à la donnée d'une section $i$ de $\mathcal{O}_{S}$ telle que}
\[
  \text{\struck{$i^{2} = -1$}}
\]
\struck{et on prend}
\[
  \text{\struck{$\omega = i\,\omega_{M}^{\otimes 2}$}}
\]
\struck{d'où}
\[
  \text{\struck{$\delta'(q_{M}) = \omega^{\otimes 2} = -\omega_{M}^{\otimes 2}$}}
\]
\struck{et on peut \ill{} \ill{} \ill{} iso}
\[
  \text{\struck{$(\gamma M, -\det_{\gamma M})$}} \quad \text{\struck{$\delta'(-q_{M}) = -\omega^{\otimes(-2)}$}}
\]
\struck{donc on a un iso $\gamma M$}
\[
  \text{\struck{$(\gamma M, -q_{M} = \det_{\gamma M}, \omega = i\omega_{M}) \longrightarrow (\gamma M', q_{M'}, \omega_{M'})$}}
\]
\struck{(quitte à \uncertain{prendre} $M'$ : iso canonique près). On peut prendre}
\[
  \text{\struck{$M' = M$,}} \qquad
  \text{\struck{$\gamma M \xrightarrow{\ \theta\ } \gamma M' = \gamma M$, \quad $x \longmapsto ix$}}
\]
Car
\[
  q_{M'}(\theta u) = q_{M}(iu) = i^{2} q_{M}(u) = -q_{M}(u),
\]
\[
  \det \theta(\omega) = i^{3}\omega = i^{4}\omega_{M} = \omega_{M} = \omega_{M'}.
\]
\note{En bas à droite, « TSVP » (tournez s'il vous plaît). Devant « Car », un symbole biffé et hachuré.}

\page{112}
\struck{Donc par définition de l'iso du présent §1', \uncertain{appliqué} à $(M, q_{M}, i\omega_{M})$, on \uncertain{l'obtient} par transport de structure.}

On doit avoir évidemment
\[
  f_{q,\omega} = \theta_{q,\omega}\, g_{q}
\]
où $\theta_{q,\omega}$ est un automorphisme de $\Sigma$, qui est compatible à tout chgt de base, et fonctoriel en $(E, q, \omega)$. Pour prouver $\theta_{q,\omega} = \mathrm{id}$, OPS que $S = \operatorname{Spec} \mathbb{Z}[\tfrac{1}{2}]$, $(E, q, \omega)$ étant la forme standard. On a une section $a_{0}$ de $\Sigma$, et $\Sigma$ s'identifie au quotient $G/T$, où $G = \mathrm{SO}(3)_{S} = \mathrm{Aut}(E, q, \omega)$, et $T$ est le stabilisateur de $a_{0}$, i.e.\ un tore maximal. Les \struck{seules} automorphismes de \add{un groupe des} $\Sigma \simeq G/T$ commutant à $G$ \add{$\mathbb{Z}/2\mathbb{Z}$} : $G \to N(T)/T \simeq (\mathbb{Z}/2\mathbb{Z})_{S}$, \uncertain{engendré} précisément par l'antipodisme.

Ceci montre que $\theta_{q,\omega}$, qui est donné par une section de $N(T)/T$ sur $S$, est soit l'identité, soit l'antipodisme. Mais lequel des deux ?
\note{« OPS » : on peut supposer. La lecture « $S = \operatorname{Spec} \mathbb{Z}[\frac{1}{2}]$ » est celle que demandent les car.\ résid.\ $\neq 2$ ; le dénominateur est petit. « $\mathbb{Z}/2\mathbb{Z}$ » est ajouté au-dessus de la ligne ; la phrase qui le porte est elliptique.}

\page{113}
Sur \uncertain{une} donnée (\uncertain{supposons extension préalable de} $S$ en $S[\mathbf{i}]$ …) $\omega = i\omega_{0}$ de sorte que $\delta'(q) = \omega^{\otimes 2}$, il faut calculer l'effet de $g_{q}$ sur $u + \mathbf{i}v$, i.e.\ de $f_{\omega}$ sur $u + iv$, $u - iv$. La vérité est la suivante : on considère
\[
  (E, -q, \omega = i\omega_{0}) \simeq (\gamma M', q_{M'}, \omega_{M'})
\]
et on transforme $u + iv$, $u - iv$ en \uncertain{éléments} de $\gamma M$…

OPS
\[
  (E, q, \omega_{0}) = (\gamma M, q_{M}, \omega_{M}),
\]
i.e.\
\[
  (E, -q, \omega) = (\gamma M, -q_{M}, i\omega_{M})
\]
\add{et on} prend $M' = M$, $\gamma M' \to \gamma M$ donné par $\theta = i\,\mathrm{id}_{\gamma M}$, qui transforme $(u + iv, u - iv)$ en
\[
  \alpha = iu - v, \qquad \beta = iu + v,
\]
\note{Dans $\beta = iu + v$, le signe est surchargé ; la suite ($2v = -\alpha + \beta$) demande $+$.}
d'où
\[
  2iu = \alpha + \beta, \qquad 2v = -\alpha + \beta,
\]
\[
  \begin{cases}
    4i[u,v] = 2[\alpha,\beta] \quad \text{i.e.} \quad 2i[u,v] = [\alpha,\beta] \\
    -4q(u) = q(\alpha) + q(\beta) + \varphi(\alpha,\beta) = \operatorname{tr} \alpha\beta \\
    4q(v) = q(\alpha) + q(\beta) - \varphi(\alpha,\beta) = -\operatorname{tr} \alpha\beta
  \end{cases}
\]
\note{Dans la deuxième ligne, $q(\alpha) + q(\beta)$ est barré d'un trait léger (ces termes sont nuls, $\alpha$, $\beta$ étant nilpotents).}
\[
  [u,v]/2q(u) = \frac{1}{2i} \cdot \frac{-2}{\operatorname{tr} \alpha\beta}\,[\alpha,\beta]
  = \frac{i}{\operatorname{tr} \alpha\beta}\,[\alpha,\beta] \quad \ldots
\]
\note{Le dénominateur du premier facteur est récrit ; on lit $2i$. Dans le second, un mot biffé devant $\alpha\beta$.}
i.e.\
\[
  \theta\bigl([u,v]/2q(u)\bigr) = i\,\frac{[u,v]}{2q(u)} = -[\alpha,\beta]/\operatorname{tr} \alpha\beta
\]

\page{114}
On trouve l'identité \uncertain{dès qu'on} trouve comme isom.\ $g_{M}$
\[
  X_{M} \times X_{M} \longrightarrow \hat{\Sigma}_{M}
\]
celui qui, pour $\alpha, \beta \in \Gamma\,\gamma M$, définissant des sections de $X$ (i.e.\ tels que $\alpha$, $\beta$ aient tr.\ et dét.\ nuls, et soient $\neq 0$ sur \uncertain{toute} fibre) \uncertain{assez} disjointes fibre par fibre (i.e.\
\[
  \operatorname{tr}(\alpha\beta) \neq 0 \qquad \text{fibre par fibre}),
\]
associe la \struck{\ill{}} \add{section} de $\Sigma \subset \gamma M$ définie par
\[
  -[\alpha,\beta]/\operatorname{tr}(\alpha\beta)
\]
(et non $+[\alpha,\beta]/\operatorname{tr}\alpha\beta$). Il faut vérifier sur une propriété \uncertain{tirée} \uncertain{conséquente} \ill{} \ill{} \ill{} l'interprétation comme section de $\hat{\Sigma}$, il faut prendre
\[
  -[\alpha,\beta]/\operatorname{tr}(\alpha\beta) + 1 \quad \text{ou encore}
\]
\[
  \operatorname{tr}(\alpha\beta).1 - [\alpha,\beta].
\]
Donc l'iso \struck{\ill{}}
\[
  X_{M} \times X_{M} \longrightarrow \hat{\Sigma}_{M}
\]
\struck{\ill{}} provient d'un \struck{iso} \add{hom.}\ \struck{\ill{}} \add{\uncertain{bilinéaire}} au niveau \struck{\ill{}} \add{des cônes} affines $C = \check{\mathbb{V}}_{\mathrm{aff}}(q_{\gamma M})$ et
\[
  D \simeq \mathbb{V}_{\mathrm{aff}}(q_{M})
\]
\struck{$\tilde{X}_{M} \times \tilde{X}_{M} \longrightarrow \ill{}$}
\note{La dernière ligne est surchargée et barrée ; la flèche biffée allait vers un symbole hachuré. L'indice « aff » sous $\mathbb{V}$ est une lecture.}

\page{115}
\[
  \text{\struck{$(\alpha,\beta) \longmapsto$}}
\]
\[
  F_{M} : C_{M} \times C_{M} \longrightarrow D_{M}
\]
\[
  F_{M}(\alpha,\beta) = \operatorname{tr}(\alpha\beta).1 - [\alpha,\beta]
\]
compatible avec l'homom.
\[
  \mathbb{G}_{m} \times \mathbb{G}_{m} \longrightarrow \mathbb{G}_{m}, \qquad (\lambda,\mu) \longmapsto \lambda\mu
\]
sur les groupes d'homothéties. \struck{\ill{}}

\struck{\uncertain{cônes} épointés, celui-ci induisant par passage aux quotients des $\mathbb{G}_{m}$-torseurs $\hat{C}_{M} \times \hat{C}_{M} / \mathbb{G}_{m}$ $\simeq \hat{D}_{M}$ (épointé, \uncertain{antidiagonale} \uncertain{exclue})}
\note{Ce passage, encadré d'un trait sinueux, est barré de plusieurs diagonales ; la lecture des mots sous les ratures est incertaine.}

Le fait que l'on s'envoie bien dans $D_{M}$ provient de la formule ($\alpha$, $\beta$ étant nilpotents, i.e.\ $\det\alpha = \det\beta = \operatorname{tr}\alpha = \operatorname{tr}\beta = 0$)
\[
  \Longrightarrow\ q_{M}\bigl([\alpha,\beta]/\varphi(\alpha,\beta)\bigr) = 1
  \qquad \text{i.e.} \qquad \det([\alpha,\beta]) + \operatorname{tr}(\alpha\beta)^{2} = 0
\]
i.e.\
\[
  q_{M}[\alpha,\beta] = \varphi(\alpha,\beta)^{2}
\]
\note{Devant $\varphi(\alpha,\beta)^{2}$, un signe $-$ est biffé.}
cas particulier de la formule générale (valable pour $\alpha, \beta \in \Gamma\,\gamma M$ \uncertain{arbitraires})
\[
  q_{M}([\alpha,\beta]) = -4\,q(\alpha)\,q(\beta) + \varphi(\alpha,\beta)^{2}
\]

\page{116}
\[
  (1) \qquad \det([\alpha,\beta]) + (\operatorname{tr}(\alpha\beta))^{2} - 4(\det\alpha)(\det\beta) = 0
\]
si $\operatorname{tr}\alpha = \operatorname{tr}\beta = 0$. Soit
\[
  u = \alpha + s.1, \qquad v = \beta + t.1 \qquad \text{donc} \qquad
  \operatorname{tr} u = 2s, \quad \operatorname{tr} v = 2t,
\]
d'où
\[
  (4) \qquad [u,v] = [\alpha,\beta]
\]
\[
  \text{\struck{$\operatorname{tr}(\alpha\beta) = \operatorname{tr}((u-s)(v-t)) = \operatorname{tr}(uv) + 2st$}}
\]
\[
  \text{\struck{$\operatorname{tr}\alpha\beta = \operatorname{tr} uv - s\operatorname{tr} v - t\operatorname{tr} u + 2st$}}
\]
\[
  \text{\struck{$= \operatorname{tr} uv - \operatorname{tr} u \operatorname{tr} v$}}
\]
\note{Les trois lignes barrées portent sous leurs termes des accolades « $\frac{1}{2}\operatorname{tr} u \operatorname{tr} v$ », « $\frac{1}{2}\operatorname{tr} v \operatorname{tr} u$ », « $\frac{1}{2}\operatorname{tr} u \operatorname{tr} v$ », barrées avec elles.}
\[
  uv = \alpha\beta + s\beta + t\alpha + st \qquad \text{donc}
\]
\[
  \operatorname{tr} uv = \operatorname{tr}\alpha\beta + 2st = \operatorname{tr}\alpha\beta + \tfrac{1}{2}\operatorname{tr} u \operatorname{tr} v
\]
i.e.\
\[
  \operatorname{tr}\alpha\beta = \operatorname{tr} uv - \tfrac{1}{2}\operatorname{tr} u \operatorname{tr} v
\]
\[
  (2) \qquad (\operatorname{tr}\alpha\beta)^{2} = (\operatorname{tr} uv)^{2} - (\operatorname{tr} uv)(\operatorname{tr} u)(\operatorname{tr} v) + \tfrac{1}{4}(\operatorname{tr} u)^{2}(\operatorname{tr} v)^{2}
\]
\[
  (3) \qquad
  \begin{cases}
    \det\alpha = \det(u-s) = \det(s-u) = s^{2} + \det u = \tfrac{1}{4}(\operatorname{tr} u)^{2} + \det u \\
    \det\beta = t^{2} + \det v = \tfrac{1}{4}(\operatorname{tr} v)^{2} + \det v
  \end{cases}
\]
\note{Les formules (3) sont transcrites telles qu'écrites ; pour une matrice $2 \times 2$ on attendrait $\det(u - s) = \det u - s^{2}$, et le signe intervient dans le résultat encadré ci-dessous.}
donc
\[
  \det([u,v]) + (\operatorname{tr} uv)^{2} - (\operatorname{tr} uv)(\operatorname{tr} u)(\operatorname{tr} v) + \tfrac{1}{4}(\operatorname{tr} u)^{2}(\operatorname{tr} v)^{2}
\]
\[
  - \tfrac{1}{4}\bigl[(\operatorname{tr} u)^{2} + 4\det u\bigr]\bigl[(\operatorname{tr} v)^{2} + 4\det v\bigr]
\]
\note{Les termes $\frac{1}{4}(\operatorname{tr} u)^{2}(\operatorname{tr} v)^{2}$ et le produit des deux $(\operatorname{tr})^{2}$ des crochets sont reliés par des accolades : ils se compensent. Le facteur $\frac{1}{4}$ devant les crochets est surchargé.}
Encadré :
\[
  \det([u,v]) - (\det u)(\operatorname{tr} v)^{2} - (\det v)(\operatorname{tr} u)^{2} - 4(\det u)(\det v)
\]
\[
  + (\operatorname{tr} uv)^{2} - (\operatorname{tr} uv)(\operatorname{tr} u)(\operatorname{tr} v) = 0
\]
\note{Le signe entre $\det([u,v])$ et $(\det u)(\operatorname{tr} v)^{2}$ est une tache d'encre ; on lit $-$.}

\section*{Quadriques affines bi-régulières}

\page{117}
\note{Titre de sa main ; il est précédé d'un numéro encadré « 1'' », biffé. « bi- » est ajouté au-dessus de « régulières », qu'il souligne ; la glose entre parenthèses est ajoutée au-dessus de la ligne.}
\struck{1''} Quadriques affines \add{bi-}régulières (i.e.\ (lisses et lisse à l'infini…) \add{projectivement})

1. Soient, sur un schéma $S$, $P$ un fibré projectif de rang $2\nu - 1$, $Q$ une quadrique lisse dans $P$, $H$ un hyperplan de $P$ non tangent à $Q$, \struck{$Q \setminus H = E$} $E$ le fibré affine correspondant. Donc $E$ (\uncertain{connu} $(P, H)$) est déterminé par une extension de $\mathcal{O}_{S}$-fibrés vect.
\[
  (1) \qquad 0 \longrightarrow V \longrightarrow T \xrightarrow{\ \pi\ } \mathcal{O}_{S} \longrightarrow 0
\]
où $V$ est un $\mathcal{O}_{S}$-Module loc.\ libre (de rang $2\nu - 1$) \add{$T$ étant de rg $2\nu$.} des translations de $T$, $E = \pi^{-1}(1)$. Connaissant $Q$, elle est déterminée par la donnée d'une section de
\[
  \check{\mathbb{P}}(\mathrm{Sym}^{2}(\check{T})) = \mathbb{P}(\struck{\ill{}}\,\Gamma^{2}(T)),
\]
i.e.\ une forme quadratique \add{$q$} sur $T$ définie loc.\ \uncertain{à} multiplication près par une scalaire inversible. La condition $Q$ lisse signifie $q$ lisse i.e.\ la forme bil.\ associée $\varphi$ non dégénérée, i.e.\ \uncertain{encore} que
\[
  (2) \qquad \delta(q) \in \Gamma(S, (\det T)^{\otimes -2})
\]
est une \struck{section} \add{base} de $(\det T)^{\otimes -2}$, i.e.\ est partout non nulle. La condition que $H$ est non tangent : $Q$ signifie aussi que $q \vert V$ est une forme lisse. \struck{On va le \ill{}} \struck{que le discriminant \ill{}}

\uncertain{Notamment}, en notant $\mathbb{D}$ le \uncertain{sous-fibré} inversible

\note{Le « rang $2\nu - 1$ » de $P$ et de $V$ est récrit (« $2\nu-1$ » surchargé). Dans (2), le $T$ de $\det T$ est surchargé ; « det » est souligné dans la ligne suivante.}

\page{118}
de $T$ défini par
\[
  (3) \qquad \mathbb{D} = V^{\perp}
\]
et on a, \uncertain{par} \ill{}
\[
  (4) \qquad \mathbb{D} \simeq (E/V)^{\vee} \simeq \mathcal{O}_{S}^{\vee} \simeq \mathcal{O}_{S}
\]
i.e.\ $\mathbb{D}$ est muni \uncertain{canoniquement} d'une \struck{section} \add{base} $a$, i.e.\ d'une section
\[
  (4') \qquad
  \begin{cases}
    a \in \Gamma(S, \mathbb{D}) \subset \Gamma(T) \\
    E = \{x \in \Gamma(-, T) \mid \varphi(x, a) = 1\}
  \end{cases}
\]
--- du moins quand $q$ est \emph{donnée} \struck{\ill{}} (\uncertain{non} seulement mod.\ \uncertain{unités}). Dans le cas général, regardant $q$ comme une \uncertain{section} \ill{} d'un sous-fibré inv.\ $L$ (quotient \ill{}) de $\Gamma^{2}(T)$, i.e.\ dual de $\mathrm{Sym}^{2}(\check{T})$…), on trouve
\[
  \mathbb{D} \otimes (E/V) \simeq L \quad \text{i.e.}
\]
[comme $T/V \simeq \mathcal{O}_{S}$] on trouve un isom.\ canonique
\[
  (5) \qquad \mathbb{D} \xrightarrow{\ \sim\ } L
\]
La condition $H$ non tangent : $Q$ s'exprime \uncertain{donc} par la condition que
\note{Dans (4), (4') et la formule $\mathbb{D} \otimes (E/V)$, le $E$ (ou $T$) est surchargé et peu sûr ; la lecture « $T/V \simeq \mathcal{O}_{S}$ » de la ligne suivante suggère $T$. La phrase se poursuit p.~119.}

\page{119}
\marginal{vérifier cette affirmation !}
$q \vert \mathbb{D}$ \uncertain{non nulle} fibre par fibre --- \uncertain{ou encore}, en identifiant $q \vert \mathbb{D}$ à un homom.
\[
  q : \mathbb{D}^{\otimes 2} \longrightarrow L,
\]
que cet homom.\ soit un \emph{iso}
\[
  (6) \qquad q : \mathbb{D}^{\otimes 2} \xrightarrow{\ \sim\ } L
\]
La condition \uncertain{E}… \uncertain{combinée avec} l'iso.\ (5) \uncertain{donne} un iso.
\[
  L^{\otimes 2} \simeq L
\]
i.e.\ un iso
\[
  (7) \qquad \mathbb{D} \simeq L \simeq \mathcal{O}_{S}
\]
Ce qui peut \uncertain{aussi} s'exprimer ainsi : \add{\uncertain{Pour} que $H = \mathbb{P}(W)$ soit non tangent : $Q \subset P$, il faut et il suffit}

\emph{Proposition} (\uncertain{qu'il} existe une \struck{unique} section $a$ de $\mathbb{D}$ \struck{\ill{}}, et une section $q$ du sous-fibré $\check{L}$ de $\mathrm{Sym}^{2}(\check{T})$ définissant $Q$, telles que
\[
  (8) \qquad \text{a)}\ q(a) = 1
\]
\[
  (9) \qquad \text{b)}\ \text{\struck{$E = \{x \in \Gamma(-,T)$}}\ \varphi(a,x) \text{\struck{$= 1\}$}}
  \quad (x \in \Gamma(-,T))
\]
Alors $a$, \add{\uncertain{où} $\varphi$ est la forme bil.\ associée à $q$,} \struck{\ill{}} \uncertain{sont} déterminés de façon unique, et \uncertain{ce sont} respectivement des bases de $\mathbb{D}$ et de $\check{L}$. Enfin, on a
\[
  (10) \qquad \pi(a) = 2
\]
Je vais : \uncertain{prouver} cette dernière relation. \struck{Or \uncertain{les} plus \ill{}, pour $x \in \Gamma(-,T)$}
\note{La marge gauche porte, en travers, une note de sa main au crayon (« \uncertain{vérifier} cette \uncertain{affirmation} ! »), et, le long de la Proposition, un trait vertical et quelques mots écrits verticalement, dont on lit seulement « $q \vert V^{\perp}$ » et « \uncertain{équivaut} » ; le reste est \ill{}. Dans (9), la description de $E$ est biffée par hachures, ne laissant que $\varphi(a,x)$ ; la forme exacte de la condition restante n'est pas lisible. Le soulignement de « Proposition » est le sien.}

\page{120}
\[
  \text{\struck{$(11) \qquad \pi(x) = \varphi(a,x)$,}}
\]
\uncertain{ce qui} s'obtient en faisant $x = a$ dans (9), et notant que $\varphi(a,a) = 2q(a) = 2$.

\emph{Scholie.} La cat.\ des $(P, Q, H)$ comme ci-dessus équivaut donc à celle des $(T, q, a)$, $T$ fibré vectoriel de rg $2\nu$, $q$ forme quadratique \uncertain{lisse} sur $T$, $a$ section de $T$, telles que $q(a) = 1$.

\emph{N.B.} a) En car.\ résiduelles $\neq 2$, on a donc que (i.e.\ $E$ est muni d'une section \uncertain{privilégiée} qui en \uncertain{fait} un \uncertain{fibré vectoriel})
\[
  T \simeq V + \mathbb{D} \simeq V + \mathcal{O}_{S} \quad \text{canoniquement, la décomposition}
\]
\uncertain{étant} orthogonale et $Q$ \uncertain{étant} défini canoniquement par une forme quadratique sur $V$, $Q \cap E$ \uncertain{l'équation} (dans $x \in \Gamma(-, E)$)
\[
  (12) \qquad q(x,t) = q_{0}(x) + t^{2}
\]
\[
  q_{0}(x) + 1 = 0 \quad \text{i.e.} \quad q_{0}(x) = -1
\]
\uncertain{D'où} en car.\ rés.\ $\neq 2$, la cat.\ des triples $(P, Q, H)$, ou des couples $(E, Q \cap E)$ ($E$ étant muni de sa structure de fibré affine…) est donc équivalente à celle des couples $(V, q_{0})$, où $V$ est un fibré vectoriel de rang $2\nu - 1$.

En car.~2, par contre, on doit avoir $\mathbb{D} \subset V$, et il n'y a pas de \uncertain{décomposition} \uncertain{canonique}, \uncertain{par ex.}\ des \uncertain{réductions} où on a \uncertain{toujours} une forme de degré $2\nu - 1$.

b) Si $T$ est de rang \uncertain{impair} $2\nu + 1$, donc $V$ de rang pair $2\nu$, les conditions que $Q$ soit
\note{Dans le Scholie, le passage « $T$ fibré vectoriel … $q(a) = 1$ » est ajouté au-dessus de la ligne, en remplacement d'un premier « bases, sur $T$, $a$ section de $T$ … » biffé ; la ligne de N.B.\ commence par un « En car.\ résid.\ $\neq 2$ » biffé puis repris. La parenthèse « (i.e.\ $E$ est muni …) » est un ajout interlinéaire. Le numéro (12) surcharge un autre chiffre ; « (11) » en tête de page est biffé avec la formule. Dans (12), l'argument de $q_{0}$ est surchargé. La phrase b) se poursuit au-delà de ce lot.}

\end{document}
